\section{Background}

This section provides essential technical background on blockchain technology and Anti-Money Laundering/Countering the Financing of Terrorism regulatory frameworks necessary for understanding Web3 RegTech systems.

\subsection{Blockchain and Web3 Fundamentals}

Blockchain technology fundamentally restructures data integrity and transaction validation through distributed consensus mechanisms. First introduced by Nakamoto \cite{nakamoto2008bitcoinpeertopeerelectronic}, blockchain maintains a continuously growing ledger of cryptographically linked blocks, each containing timestamped transaction data and a reference to its predecessor. This architecture creates an immutable record that can be independently verified by any network participant, eliminating reliance on centralized authorities for transaction validation.

Web3 represents the evolution of internet infrastructure toward decentralized architectures built atop blockchain foundations \cite{wood2014ethereumsecuredecentralised}. Smart contracts constitute the primary programmability layer, enabling self-executing code that automatically enforces agreement terms without intermediary oversight \cite{szabo1997formalizingsecuringrelationships}. Decentralized Finance protocols leverage smart contracts to recreate traditional financial services including lending markets, trading venues, and derivative instruments entirely through algorithmic intermediation \cite{harvey2021defifuturefinance}. Decentralized Autonomous Organizations encode governance structures in smart contracts, enabling collective decision-making through token-weighted voting mechanisms \cite{voshmgir2020tokeneconomyhow}. Non-Fungible Tokens provide cryptographically secured ownership records for unique digital or physical assets \cite{wang2021nonfungibletokennft}.

The transparency inherent in public blockchains fundamentally differs from traditional financial systems. Every transaction broadcasts to all network participants, creating comprehensive visibility into fund flows and contract interactions \cite{narayanan2016bitcoincryptocurrencytechnologies}. This transparency enables novel analytical approaches while simultaneously creating privacy challenges, as the permanent public record of all activities can reveal sensitive information when addresses become associated with real-world identities.

\subsection{AML/CFT Regulatory Framework}

Anti-Money Laundering and Countering the Financing of Terrorism regulations establish baseline obligations that financial institutions must satisfy to prevent illicit use of financial systems \cite{financialactiontaskforcefatf2012internationalstandardscombating}. Customer Due Diligence requires verification of customer identities and understanding the nature of business relationships, enabling institutions to assess appropriate risk levels and design proportionate monitoring strategies. Transaction Monitoring demands continuous surveillance of customer activities to identify patterns inconsistent with known legitimate business or personal activities. Suspicious Activity Reporting mandates that institutions notify authorities when transactions potentially indicate money laundering, terrorist financing, fraud, or other criminal activities. Record Keeping obligations require maintenance of identification data and transaction records for specified periods, typically five years, ensuring availability for retrospective investigation.

The Financial Action Task Force, an intergovernmental organization, establishes global standards for AML/CFT compliance \cite{financialactiontaskforcefatf2012internationalstandardscombating}. FATF's Recommendation 16, commonly known as the Travel Rule, requires Virtual Asset Service Providers to obtain, hold, and exchange information about originators and beneficiaries of virtual asset transfers \cite{force2021updatedguidanceriskbased}. Specifically, VASPs must transmit originator's name, account number or address, and beneficiary's name and account number for transfers exceeding specified thresholds. This requirement attempts to replicate wire transfer transparency in cryptocurrency contexts, though technical implementation proves significantly more challenging due to decentralized network architectures and interoperability gaps among Travel Rule protocols \cite{force2021opportunitieschallengesnew}.

Virtual Asset Service Providers, as defined by FATF, are entities conducting activities or operations on behalf of customers including exchange between virtual and fiat currencies, transfer of virtual assets, safekeeping or administration of virtual assets, and participation in financial services related to virtual assets \cite{force2021updatedguidanceriskbased}. VASPs face compliance obligations similar to traditional financial institutions but must address unique challenges posed by blockchain technology including pseudonymous addresses lacking inherent identity information, cross-chain transactions challenging single-ledger monitoring approaches, self-custody eliminating VASP visibility for off-platform activities, and smart contract complexity creating interpretation challenges.

\subsection{Compliance Challenges in Web3}

The architectural characteristics of Web3 systems create fundamental tensions between regulatory compliance requirements and technological design principles. Privacy-enhancing technologies provide legitimate financial privacy protections but simultaneously enable obfuscation of illicit fund flows. Coin mixing services, privacy-focused cryptocurrencies like Monero and Zcash, and zero-knowledge proof systems all complicate transaction tracing and risk assessment \cite{moser2018empiricalanalysistraceability,kappos2018empiricalanalysisanonymity}. Recent research explores privacy-preserving compliance approaches balancing these competing interests through selective disclosure mechanisms, decentralized identifiers, and zero-knowledge proofs of compliance \cite{yang2024researchidentitydata,ben-sasson2017scalablezeroknowledge,zhang2024blockchainzeroknowledge}.

Decentralization challenges traditional accountability frameworks that assume identifiable intermediaries responsible for compliance. Decentralized Finance protocols execute autonomously through smart contracts without ongoing human intervention, raising questions about who bears compliance responsibility \cite{harvey2021defifuturefinance,werbach2018trustverifywhy}. Protocol developers may have limited control after deployment, liquidity providers supply capital without operational control, and users interact directly with contracts without intermediary oversight. This disintermediation creates regulatory gaps where no clearly accountable party exists to enforce compliance requirements.

Technological innovation proceeds at a pace that regulatory frameworks and compliance tools struggle to match. New token standards, novel DeFi mechanisms, Layer-2 scaling solutions, and cross-chain bridges continuously introduce unprecedented interaction patterns requiring compliance interpretation \cite{mohammedabdul2024navigatingblockchainstwin,xu2024exploringblockchaintechnology}. Each innovation potentially introduces new risks or evasion vectors that existing detection systems fail to recognize, necessitating continual model retraining and intelligence sharing.
